% \begin{document}
\section{A simple imperative language IMP}

We will now consider a more realistic programming language, one where we can assign values to
variables and execute control constructs such as \textbf{if} and \textbf{while}. The syntax for
this imperative language, called IMP, is as follows:

\begin{align*}
\text{arithmetic expressions} \quad & a \in \mathbf{Aexp} \quad a ::= x \mid n \mid a_1 + a_2
\mid a_1 \times a_2 \\
\text{Boolean expressions} \quad & b \in \mathbf{Bexp} \quad b ::= \mathbf{true} \mid
\mathbf{false} \mid a_1 < a_2 \\
\text{commands} \quad & c \in \mathbf{Com} \quad c ::= \mathbf{skip} \mid x := a \mid c_1; c_2
\mid \mathbf{if}\ b\ \mathbf{then}\ c_1\ \mathbf{else}\ c_2 \mid \mathbf{while}\ b\
\mathbf{do}\ c
\end{align*}

\subsection{Small-step operational semantics}

We'll first give a small-step operational semantics for IMP. The configurations in this
language are of the form $\langle \sigma, c \rangle$, $\langle \sigma, b \rangle$, and $\langle
\sigma, a \rangle$, where $\sigma$ is a store. The final configurations are of the form
$\langle \sigma, \mathbf{skip} \rangle$ for commands, $\langle \sigma, \mathbf{true} \rangle$
and $\langle \sigma, \mathbf{false} \rangle$ for Boolean expressions, and $\langle \sigma, n
\rangle$ for arithmetic expressions. There are three different small-step operational
semantics relations: one each of the syntactic categories.
\[
\to_{\mathbf{Com}} \; \subseteq \; (\mathbf{Store} \times \mathbf{Com}) \times (\mathbf{Store}
\times \mathbf{Com})
\]
\[
\to_{\mathbf{Bexp}} \; \subseteq \; (\mathbf{Store} \times \mathbf{Bexp}) \times
(\mathbf{Store} \times \mathbf{Bexp})
\]
\[
\to_{\mathbf{Aexp}} \; \subseteq \; (\mathbf{Store} \times \mathbf{Aexp}) \times
(\mathbf{Store} \times \mathbf{Aexp})
\]

For brevity, we will overload the symbol $\to$ and use it to refer to all of these relations.
Which relation is being used will be clear from context. The evaluation rules for arithmetic
and Boolean expressions are similar to the ones we've seen before. However, note that since the
arithmetic expressions no longer contain assignment, arithmetic and Boolean expressions can not
update the store. But to streamline everything we keep it as an output. In the small-step we'll fix the fact that the store doesn't change, but this can be altered if needed using $ \sigma' $.

\medskip
\noindent\textbf{Arithmetic expressions}

\[
\frac{n = \sigma(x)}{\langle \sigma, x \rangle \to \langle \sigma, n \rangle}
\]

\[
\frac{\langle \sigma, a_1 \rangle \to \langle \sigma, a_1' \rangle}{\langle \sigma, a_1 + a_2
\rangle \to \langle \sigma, a_1' + a_2 \rangle}
\qquad
\frac{\langle \sigma, a_2 \rangle \to \langle \sigma, a_2' \rangle}{\langle \sigma, n + a_2
\rangle \to \langle \sigma, n + a_2' \rangle}
\qquad
\frac{p = n + m}{\langle \sigma, n + m \rangle \to \langle \sigma, p \rangle}
\]

\[
\frac{\langle \sigma, a_1 \rangle \to \langle \sigma, a_1' \rangle}{\langle \sigma, a_1 \times
a_2 \rangle \to \langle \sigma, a_1' \times a_2 \rangle}
\qquad
\frac{\langle \sigma, a_2 \rangle \to \langle \sigma, a_2' \rangle}{\langle \sigma, n \times a_2
\rangle \to \langle \sigma, n \times a_2' \rangle}
\qquad
\frac{p = n \times m}{\langle \sigma, n \times m \rangle \to \langle \sigma, p \rangle}
\]

\medskip
\noindent\textbf{Boolean expressions}

\[
\frac{\langle \sigma, a_1 \rangle \to \langle \sigma, a_1' \rangle}{\langle \sigma, a_1 < a_2
\rangle \to \langle \sigma, a_1' < a_2 \rangle}
\qquad
\frac{\langle \sigma, a_2 \rangle \to \langle \sigma, a_2' \rangle}{\langle \sigma, n < a_2
\rangle \to \langle \sigma, n < a_2' \rangle}
\]

\[
\frac{n < m}{\langle \sigma, n < m \rangle \to \langle \sigma, \mathbf{true} \rangle}
\qquad
\frac{n \geq m}{\langle \sigma, n < m \rangle \to \langle \sigma, \mathbf{false} \rangle}
\]

\medskip
\noindent\textbf{Commands}

\[
\frac{\langle \sigma, a \rangle \to \langle \sigma, a' \rangle}{\langle \sigma, x := a \rangle
\to \langle \sigma, x := a' \rangle}
\qquad
\frac{}{\langle \sigma, x := n \rangle \to \langle \sigma[x \mapsto n], \mathbf{skip} \rangle}
\]

\[
\frac{\langle \sigma, c_1 \rangle \to \langle \sigma', c_1' \rangle}{\langle \sigma, c_1; c_2
\rangle \to \langle \sigma', c_1'; c_2 \rangle}
\qquad
\frac{}{\langle \sigma, \mathbf{skip}; c_2 \rangle \to \langle \sigma, c_2 \rangle}
\]

For \textbf{if} commands, we reduce the test until we get $\mathbf{true}$ or $\mathbf{false}$
and then we execute the appropriate branch:

\[
\frac{\langle \sigma, b \rangle \to \langle \sigma, b' \rangle}{\langle \sigma, \mathbf{if}\ b\
\mathbf{then}\ c_1\ \mathbf{else}\ c_2 \rangle \to \langle \sigma, \mathbf{if}\ b'\
\mathbf{then}\ c_1\ \mathbf{else}\ c_2 \rangle}
\]

\[
\frac{}{\langle \sigma, \mathbf{if}\ \mathbf{true}\ \mathbf{then}\ c_1\ \mathbf{else}\ c_2
\rangle \to \langle \sigma, c_1 \rangle}
\qquad
\frac{}{\langle \sigma, \mathbf{if}\ \mathbf{false}\ \mathbf{then}\ c_1\ \mathbf{else}\ c_2
\rangle \to \langle \sigma, c_2 \rangle}
\]

For \textbf{while} loops, the above strategy doesn't work, because if we evaluate $ b $ we won't be able to re-evaluate it for the next loop. Instead, we use the following
rule, which can be thought of as ``unrolling'' the loop, one iteration at a time.

\[
\frac{}{\langle \sigma, \mathbf{while}\ b\ \mathbf{do}\ c \rangle \to \langle \sigma,
\mathbf{if}\ b\ \mathbf{then}\ (c; \mathbf{while}\ b\ \mathbf{do}\ c)\ \mathbf{else}\
\mathbf{skip} \rangle}
\]
\nt{
With this new rule, induction on the structure of the term doesn't work! This is because instead of always getting smaller, the step is to a bigger $ c $. Sometimes this is managable, but not always.
}

We can now take a concrete program and see how it executes under the above rules. Consider we
execute the program

\ex{Small-step eval of IMP}{
\[
\mathit{foo} := 3;\ \mathbf{while}\ \mathit{foo} < 4\ \mathbf{do}\ \mathit{foo} := \mathit{foo}
+ 5
\]
The execution works as follows:
\begin{align*}
&\langle \sigma, \mathit{foo} := 3; \mathbf{while}\ \mathit{foo} < 4\ \mathbf{do}\ \mathit{foo}
:= \mathit{foo} + 5 \rangle \\
\to\ &\langle \sigma', \mathbf{skip}; \mathbf{while}\ \mathit{foo} < 4\ \mathbf{do}\ \mathit{foo}
:= \mathit{foo} + 5 \rangle && \text{where } \sigma' = \sigma[\mathit{foo} \mapsto 3] \\
\to\ &\langle \sigma', \mathbf{while}\ \mathit{foo} < 4\ \mathbf{do}\ \mathit{foo} := \mathit{foo}
+ 5 \rangle \\
\to\ &\langle \sigma', \mathbf{if}\ \mathit{foo} < 4\ \mathbf{then}\ (\mathit{foo} := \mathit{foo}
+ 5; W)\ \mathbf{else}\ \mathbf{skip} \rangle \\
\to\ &\langle \sigma', \mathbf{if}\ 3 < 4\ \mathbf{then}\ (\mathit{foo} := \mathit{foo} + 5;
W)\ \mathbf{else}\ \mathbf{skip} \rangle \\
\to\ &\langle \sigma', \mathbf{if}\ \mathbf{true}\ \mathbf{then}\ (\mathit{foo} := \mathit{foo}
+ 5; W)\ \mathbf{else}\ \mathbf{skip} \rangle \\
\to\ &\langle \sigma', \mathit{foo} := \mathit{foo} + 5; \mathbf{while}\ \mathit{foo} < 4\
\mathbf{do}\ \mathit{foo} := \mathit{foo} + 5 \rangle \\
\to\ &\langle \sigma', \mathit{foo} := 3 + 5; \mathbf{while}\ \mathit{foo} < 4\ \mathbf{do}\
\mathit{foo} := \mathit{foo} + 5 \rangle \\
\to\ &\langle \sigma', \mathit{foo} := 8; \mathbf{while}\ \mathit{foo} < 4\ \mathbf{do}\
\mathit{foo} := \mathit{foo} + 5 \rangle \\
\to\ &\langle \sigma'', \mathbf{skip}; \mathbf{while}\ \mathit{foo} < 4\ \mathbf{do}\
\mathit{foo} := \mathit{foo} + 5 \rangle && \text{where } \sigma'' = \sigma'[\mathit{foo}
\mapsto 8] \\
\to\ &\langle \sigma'', \mathbf{while}\ \mathit{foo} < 4\ \mathbf{do}\ \mathit{foo} :=
\mathit{foo} + 5 \rangle \\
\to\ &\langle \sigma'', \mathbf{if}\ \mathit{foo} < 4\ \mathbf{then}\ (\mathit{foo} :=
\mathit{foo} + 5; W)\ \mathbf{else}\ \mathbf{skip} \rangle \\
\to\ &\langle \sigma'', \mathbf{if}\ 8 < 4\ \mathbf{then}\ (\mathit{foo} := \mathit{foo} + 5;
W)\ \mathbf{else}\ \mathbf{skip} \rangle \\
\to\ &\langle \sigma'', \mathbf{if}\ \mathbf{false}\ \mathbf{then}\ (\mathit{foo} :=
\mathit{foo} + 5; W)\ \mathbf{else}\ \mathbf{skip} \rangle \\
\to\ &\langle \sigma'', \mathbf{skip} \rangle
\end{align*}
where $W$ is an abbreviation for the while loop $\mathbf{while}\ \mathit{foo} < 4\
\mathbf{do}\ \mathit{foo} := \mathit{foo} + 5$.
}

\section{Large-step operational semantics for IMP}

We define large-step evaluation relations for arithmetic expressions, Boolean expressions, and
commands. The relation for arithmetic expressions relates an arithmetic expression and store to
the integer value that the expression evaluates to. For Boolean expressions, the final value is
in $\mathbf{Bool} = \{\mathbf{true}, \mathbf{false}\}$. For commands, the final value is a
store.

\[
\Downarrow_{\mathbf{Aexp}} \; \subseteq \; (\mathbf{Aexp} \times \mathbf{Store}) \times
\mathbf{Int}
\]
\[
\Downarrow_{\mathbf{Bexp}} \; \subseteq \; (\mathbf{Bexp} \times \mathbf{Store}) \times
\mathbf{Bool}
\]
\[
\Downarrow_{\mathbf{Com}} \; \subseteq \; (\mathbf{Com} \times \mathbf{Store}) \times
\mathbf{Store}
\]

Again, we overload the symbol $\Downarrow$ and use it for any of these three relations; which
relation is intended will be clear from context. We also use infix notation, for example
writing $\langle \sigma, c \rangle \Downarrow \sigma'$ if $(\langle \sigma, c \rangle, \sigma')
\in {\Downarrow_{\mathbf{Com}}}$.

\nt{
  In this course, we're not interested in concrete syntax, so ambiguous AST aren't a problem. Braces can be used to explicit.
}

\medskip
\noindent\textbf{Arithmetic expressions.}

\[
\frac{}{\langle \sigma, n \rangle \Downarrow n}
\qquad
\frac{\sigma(x) = n}{\langle \sigma, x \rangle \Downarrow n}
\]

\[
\frac{\langle \sigma, a_1 \rangle \Downarrow n_1 \qquad \langle \sigma, a_2 \rangle \Downarrow
n_2 \qquad n = n_1 + n_2}{\langle \sigma, a_1 + a_2 \rangle \Downarrow n}
\qquad
\frac{\langle \sigma, a_1 \rangle \Downarrow n_1 \qquad \langle \sigma, a_2 \rangle \Downarrow
n_2 \qquad n = n_1 \times n_2}{\langle \sigma, a_1 \times a_2 \rangle \Downarrow n}
\]

\medskip
\noindent\textbf{Boolean expressions.}

\[
\frac{}{\langle \sigma, \mathbf{true} \rangle \Downarrow \mathbf{true}}
\qquad
\frac{}{\langle \sigma, \mathbf{false} \rangle \Downarrow \mathbf{false}}
\]

\[
\frac{\langle \sigma, a_1 \rangle \Downarrow n_1 \qquad \langle \sigma, a_2 \rangle \Downarrow
n_2 \qquad n_1 < n_2}{\langle \sigma, a_1 < a_2 \rangle \Downarrow \mathbf{true}}
\qquad
\frac{\langle \sigma, a_1 \rangle \Downarrow n_1 \qquad \langle \sigma, a_2 \rangle \Downarrow
n_2 \qquad n_1 \geq n_2}{\langle \sigma, a_1 < a_2 \rangle \Downarrow \mathbf{false}}
\]

\medskip
\noindent\textbf{Commands.}

\[
\frac{}{\langle \sigma, \mathbf{skip} \rangle \Downarrow \sigma} \; \text{SKIP}
\qquad
\frac{\langle \sigma, a \rangle \Downarrow n}{\langle \sigma, x := a \rangle \Downarrow
\sigma[x \mapsto n]} \; \text{ASSGN}
\]

\[
\frac{\langle \sigma, c_1 \rangle \Downarrow \sigma' \qquad \langle \sigma', c_2 \rangle
\Downarrow \sigma''}{\langle \sigma, c_1; c_2 \rangle \Downarrow \sigma''} \; \text{SEQ}
\]

\[
\frac{\langle \sigma, b \rangle \Downarrow \mathbf{true} \qquad \langle \sigma, c_1 \rangle
\Downarrow \sigma'}{\langle \sigma, \mathbf{if}\ b\ \mathbf{then}\ c_1\ \mathbf{else}\ c_2
\rangle \Downarrow \sigma'} \; \text{IF-T}
\qquad
\frac{\langle \sigma, b \rangle \Downarrow \mathbf{false} \qquad \langle \sigma, c_2 \rangle
\Downarrow \sigma'}{\langle \sigma, \mathbf{if}\ b\ \mathbf{then}\ c_1\ \mathbf{else}\ c_2
\rangle \Downarrow \sigma'} \; \text{IF-F}
\]

\[
\frac{\langle \sigma, b \rangle \Downarrow \mathbf{false}}{\langle \sigma, \mathbf{while}\ b\
\mathbf{do}\ c \rangle \Downarrow \sigma} \; \text{WHILE-F}
\]

\[
\frac{\langle \sigma, b \rangle \Downarrow \mathbf{true} \qquad \langle \sigma, c \rangle
\Downarrow \sigma' \qquad \langle \sigma', \mathbf{while}\ b\ \mathbf{do}\ c \rangle \Downarrow
\sigma''}{\langle \sigma, \mathbf{while}\ b\ \mathbf{do}\ c \rangle \Downarrow \sigma''} \;
\text{WHILE-T}
\]

\nt{
  The big-step semantics treat termination differently to small-step. To prove these two are equal, we could assume termination.
}

It's interesting to see that the rule for while loops does not rely on using an if command (as
in the case of small-step semantics). Why does this rule work?

\subsection{Command equivalence}

The small-step operational semantics suggests that the loop $\mathbf{while}\ b\ \mathbf{do}\ c$
should be equivalent to the command $\mathbf{if}\ b\ \mathbf{then}\ (c; \mathbf{while}\ b\
\mathbf{do}\ c)\ \mathbf{else}\ \mathbf{skip}$. Can we show that this indeed the case that the
language is defined using the above large-step evaluation?

First, we need to to be more precise about what ``equivalent commands'' mean. Our formal model
allows us to define this concept using large-step evaluations as follows. (One can write a
similar definition using $\to^*$ in small-step semantics.)

\dfn{Equivalence of commands}{
Two commands $c$ and $c'$ are equivalent (written $c \sim c'$) if, for any stores $\sigma$ and
$\sigma'$, we have
\[
\langle \sigma, c \rangle \Downarrow \sigma' \iff \langle \sigma, c' \rangle \Downarrow
\sigma'.
\]
}

We can now state and prove the claim that $\mathbf{while}\ b\ \mathbf{do}\ c$ and $\mathbf{if}\
b\ \mathbf{then}\ (c; \mathbf{while}\ b\ \mathbf{do}\ c)\ \mathbf{else}\ \mathbf{skip}$ are
equivalent.

\thm{}{
For all $b \in \mathbf{Bexp}$ and $c \in \mathbf{Com}$ we have
\[
\mathbf{while}\ b\ \mathbf{do}\ c \; \sim \; \mathbf{if}\ b\ \mathbf{then}\ (c;
\mathbf{while}\ b\ \mathbf{do}\ c)\ \mathbf{else}\ \mathbf{skip}.
\]
}

\pf{}{
Let $W$ be an abbreviation for $\mathbf{while}\ b\ \mathbf{do}\ c$. We want to show that for
all stores $\sigma$, $\sigma'$, we have:
\[
\langle \sigma, W \rangle \Downarrow \sigma' \quad \text{if and only if} \quad \langle \sigma,
\mathbf{if}\ b\ \mathbf{then}\ (c; W)\ \mathbf{else}\ \mathbf{skip} \rangle \Downarrow \sigma'
\]
For this, we must show that both directions ($\implies$ and $\impliedby$) hold. We'll show
only direction $\implies$; the other is similar.

Assume that $\sigma$ and $\sigma'$ are stores such that $\langle \sigma, W \rangle \Downarrow
\sigma'$. It means that there is some derivation that proves this fact. Inspecting the
evaluation rules, we see that there are two possible rules whose conclusions match this fact:
\textsc{While-F} and \textsc{While-T}. We analyze each of them in turn.

\medskip
\noindent\textbf{\textsc{While-F}.} The derivation must look like the following.
\[
\cfrac{\overset{\displaystyle \vdots_1}{\langle \sigma, b \rangle \Downarrow \mathbf{false}}}
{\langle \sigma, W \rangle \Downarrow \sigma} \; \text{WHILE-F}
\]
Here, we use $\vdots_1$ to refer to the derivation of $\langle \sigma, b \rangle \Downarrow
\mathbf{false}$. Note that in this case, $\sigma' = \sigma$.

We can use $\vdots_1$ to derive a proof tree showing that the evaluation of $\mathbf{if}\ b\
\mathbf{then}\ (c; W)\ \mathbf{else}\ \mathbf{skip}$ yields the same final state $\sigma$:
\[
\cfrac{
\overset{\displaystyle \vdots_1}{\langle \sigma, b \rangle \Downarrow \mathbf{false}}
\qquad
\cfrac{}{\langle \sigma, \mathbf{skip} \rangle \Downarrow \sigma} \; \text{SKIP}
}{\langle \sigma, \mathbf{if}\ b\ \mathbf{then}\ (c; W)\ \mathbf{else}\ \mathbf{skip} \rangle
\Downarrow \sigma} \; \text{IF-F}
\]

\medskip
\noindent\textbf{\textsc{While-T}.} In this case, the derivation has the following form.
\[
\cfrac{
\overset{\displaystyle \vdots_2}{\langle \sigma, b \rangle \Downarrow \mathbf{true}}
\qquad
\overset{\displaystyle \vdots_3}{\langle \sigma, c \rangle \Downarrow \sigma''}
\qquad
\overset{\displaystyle \vdots_4}{\langle \sigma'', W \rangle \Downarrow \sigma'}
}{\langle \sigma, W \rangle \Downarrow \sigma'} \; \text{WHILE-T}
\]
We can use subderivations $\vdots_2$, $\vdots_3$, and $\vdots_4$ to show that the evaluation of
$\mathbf{if}\ b\ \mathbf{then}\ (c; W)\ \mathbf{else}\ \mathbf{skip}$ yields the same final
state $\sigma'$.
\[
\cfrac{
\overset{\displaystyle \vdots_2}{\langle \sigma, b \rangle \Downarrow \mathbf{true}}
\qquad
\cfrac{
\overset{\displaystyle \vdots_3}{\langle \sigma, c \rangle \Downarrow \sigma''}
\qquad
\overset{\displaystyle \vdots_4}{\langle \sigma'', W \rangle \Downarrow \sigma'}
}{\langle \sigma, c; W \rangle \Downarrow \sigma'} \; \text{SEQ}
}{\langle \sigma, \mathbf{if}\ b\ \mathbf{then}\ (c; W)\ \mathbf{else}\ \mathbf{skip} \rangle
\Downarrow \sigma'} \; \text{IF-T}
\]

Hence, we showed that in each of the two possible cases, the command $\mathbf{if}\ b\
\mathbf{then}\ (c; W)\ \mathbf{else}\ \mathbf{skip}$ evaluates to the same final state as the
command $W$.
}

Some questions:
\begin{itemize}
\item Can we write a program that doesn't terminate? yes, obv
\item Is IMP Turing Complete? Kinda, we don't have a way of creating new state (it's statically allocated), so if we replace the unbound \textbf{Int} type with something like \textbf{Int64}, it wouldn't be complete
\end{itemize}

% \end{document}
